\chapter{Introducción}

\section{Objetivos de la práctica}
El objetivo fundamental de esta práctica es la obtenención de información y 
familiarización con las amenazas a la seguridad digital.

\section{Fundamentos teóricos}
Para la realización de esta práctica, se aplican diversos conceptos teóricos clave 
de la ciberseguridad y la seguridad de la información, entre los que destacan:
\begin{itemize}
    \item \textbf{Conocimiento de los CVE/CWE:} CVE (Common Vulnerabilities and Exposures) es un sistema de identificación de vulnerabilidades conocido a nivel mundial, mientras que CWE (Common Weakness Enumeration) clasifica las debilidades de software. Familiarizarse con estas bases de datos es esencial para entender las amenazas actuales y cómo mitigarlas.
    \item \textbf{Phishing:} Técnica de ingeniería social utilizada para engañar a los usuarios y obtener información confidencial, como contraseñas o datos bancarios, a través de correos electrónicos, mensajes o sitios web falsos.
    \item \textbf{Malware:} Software malicioso diseñado para dañar, explotar o comprometer sistemas informáticos. Puede incluir virus, troyanos, ransomware, spyware, entre otros.
    \item \textbf{Ransomware:} Tipo de malware que cifra los archivos de la víctima y exige un rescate para restaurar el acceso a los datos. Es una amenaza creciente para empresas y particulares.
    \item \textbf{Botnets:} Redes de dispositivos comprometidos (ordenadores zombies) controlados por un atacante para realizar actividades maliciosas, como ataques DDoS, envío de spam o minería de criptomonedas.
    \item \textbf{SQL Injection:} Vulnerabilidad que permite a un atacante ejecutar código SQL malicioso en una base de datos a través de una aplicación web, lo que puede resultar en la exposición o manipulación de datos sensibles.
    \item \textbf{Cross-Site Scripting (XSS):} Vulnerabilidad que permite a un atacante inyectar scripts maliciosos en páginas web vistas por otros usuarios, lo que puede llevar al robo de cookies, secuest
\end{itemize}

\subsection{Enunciado:}
1. Introducción a CVE, CVSS y las principales bases de datos

\subsection{Introducción a CVE, CVSS y las principales bases de datos}
El estándar CVE (Common Vulnerabilities and Exposures) consiste en un sistema de
 identificación y catalogación de vulnerabilidades de software y hardware. 
 El objetivo de esta norma es suministrar un marco común para que las organizaciones, 
 desarrolladores y profesionales de la seguridad puedan discutir, reportar y 
 gestionar estas vulnerabilidades de manera uniforme. La gestión de CVE es coordinada 
 principalmente por MITRE (Management of Information Technology Research and Evaluation), 
 que establece los estándares y el marco operativo para la asignación de estos identificadores.

Por otra parte, las CNA (CVE Numbering Authorities) son organizaciones autorizadas para 
asignar identificadores CVE a vulnerabilidades específicas. Estas autoridades pueden 
ser fabricantes de software, investigadores de seguridad o entidades académicas, 
y su función es garantizar que cada vulnerabilidad se registre de manera precisa y 
rastreable. La asignación de un CVE a una vulnerabilidad implica un proceso riguroso 
de verificación y documentación para asegurar la integridad del sistema. Una vez que 
una vulnerabilidad ha sido identificada y registrada con un CVE, su impacto se 
cuantifica utilizando el sistema CVSS (Common Vulnerability Scoring System). CVSS es 
un método estandarizado que asigna una puntuación numérica (de 0.0 a 10.0) que refleja 
la severidad y la explotabilidad de la vulnerabilidad.

En cuanto a las principales bases de datos, existen diferencias clave: CVE.org es el
 registro central donde se publican y gestionan los identificadores CVE. La NVD 
 (National Vulnerability Database), mantenida por el NIST (Instituto Nacional de 
 Estándares y Tecnología), es una base de datos que recopila datos sobre 
 vulnerabilidades, incluyendo puntuaciones CVSS asociadas a múltiples identificadores 
 CVEs. Finalmente, FIRST.org (Forum of Incident Response and Security Teams) es un 
 foro más amplio centrado en la colaboración comunitaria, las mejores prácticas y el
  desarrollo de estándares de respuesta a incidentes de seguridad, aunque no es una 
  base de datos directa de vulnerabilidades como la NVD.

\section{Enunciado de la práctica}
Esta práctica consiste en la realización de un informe basado en una selección de noticias de prensa digital, posteriores al 1 de enero de 2025, relacionadas con una serie de amenazas informáticas.

Las amenazas objeto de esta actividad serán:
\begin{enumerate}
    \item Phishing, estafas a través de la red
    \item Malware difundido en redes sociales
    \item Ataques de ransomware contra corporaciones y gobiernos
    \item Botnets, ordenadores zombies
    \item Malware en dispositivos móviles
    \item Inyección de código SQL
    \item Cross-Site Scripting (XSS)
\end{enumerate}

